\paragraph{属 \(49\) 的 \texorpdfstring{$S_4$}{S4}--Hurwitz 商塔上的 Prym Weil 型、局部 \texorpdfstring{$\zeta$}{zeta} 分解与 Hodge 类传递}
沿用定理 \ref{thm:killo-s4-even-subgroup-free-action-unramified-a4-lift}、\ref{thm:killo-s4-burnside-kani-rosen-prym-square}、\ref{thm:killo-s4-genus49-jacobian-computable-isogeny} 与命题 \ref{prop:killo-two-prym-polarization-types} 的记号，记
\[
X:=\mathcal X,\qquad
H:=\mathcal H=X/A_4,\qquad
C_6:=\mathcal C_6=X/V_4,
\]
\[
Z:=X/D_8,\qquad
C_F:=\mathcal C_F=X/S_3,\qquad
Y:=X/C_4,
\]
并置
\[
A:=\operatorname{Prym}(C_6/H),\qquad
P:=\operatorname{Prym}(Y/Z).
\]

\begin{conclusion}[\texorpdfstring{$\operatorname{Prym}(C_6/H)$}{Prym(C6/H)} 的 \texorpdfstring{$\QQ(\zeta_3)$}{Q(zeta3)}-Weil 型结构与签名 \texorpdfstring{$(4,4)$}{(4,4)}]\label{con:xi-time-part9n-prym-c6h-weil-type-signature}
阿贝尔簇 \(A\) 是以
\[
K:=\QQ(\zeta_3)
\]
为虚二次域的 Weil 型阿贝尔八维簇，其 \(K\)-signature 为
\[
(4,4).
\]
并且 \(A\) 的诱导极化型为
\[
(1,1,1,1,3,3,3,3),
\]
且有
\[
\QQ(\zeta_3)\subset \operatorname{End}^0(A),\qquad
A\sim_{\QQ}J(Z)^2.
\]
特别地，存在自然二维 \(\QQ\)-向量空间
\[
W_K:=\bigwedge\nolimits_{K}^{8}H^1(A,\QQ)\subset H^8(A,\QQ)\cap H^{4,4}(A),
\]
其由 Weil classes 给出；经同源 \(A\sim_{\QQ}J(Z)^2\)，该二维 Hodge 类空间亦传递到
\[
H^8\!\bigl(J(Z)^2,\QQ\bigr)\cap H^{4,4}\!\bigl(J(Z)^2\bigr).
\]
\end{conclusion}
\begin{proof}
由定理 \ref{thm:killo-s4-even-subgroup-free-action-unramified-a4-lift}，覆盖
\[
\pi:C_6\to H
\]
为无分歧循环三重覆盖，且
\[
g(H)=5,\qquad g(C_6)=13.
\]
故
\[
\dim A=g(C_6)-g(H)=8.
\]
再由定理 \ref{thm:killo-s4-burnside-kani-rosen-prym-square} 与命题 \ref{prop:killo-two-prym-polarization-types}，
\[
A\sim_{\QQ}J(Z)^2,
\]
并且 \(A\) 的诱导极化型为 \((1^4,3^4)\)，同时
\[
\QQ(\zeta_3)\subset \operatorname{End}^0(A).
\]

为确定 signature，记 \(\sigma\) 为 \(\pi\) 的甲板变换生成元。由于 \(\pi\) 为循环三重无分歧覆盖，存在非常值 \(3\)-扭线丛
\[
L\in \operatorname{Pic}^0(H),\qquad L^{\otimes 3}\simeq \mathcal O_H,
\]
使得
\[
\pi_*\mathcal O_{C_6}\simeq \mathcal O_H\oplus L^{-1}\oplus L^{-2}.
\]
由投影公式，
\[
\pi_*\omega_{C_6}
\simeq
\omega_H\oplus (\omega_H\otimes L)\oplus (\omega_H\otimes L^2).
\]
于是 \(H^0(C_6,\Omega_{C_6}^1)\) 按 \(C_3=\langle \sigma\rangle\) 的特征分解为
\[
H^0(C_6,\Omega_{C_6}^1)=U_0\oplus U_{\chi}\oplus U_{\bar\chi},
\qquad
\chi(\sigma)=\zeta_3,
\]
其中
\[
U_0\simeq H^0(H,\omega_H),\qquad
U_{\chi}\simeq H^0(H,\omega_H\otimes L),\qquad
U_{\bar\chi}\simeq H^0(H,\omega_H\otimes L^2).
\]
由于 \(L\) 与 \(L^2\) 均为非平凡零次数线丛，Riemann--Roch 与 Serre 对偶给出
\[
h^0(H,\omega_H\otimes L)=h^0(H,\omega_H\otimes L^2)=g(H)-1=4,
\]
而
\[
h^0(H,\omega_H)=g(H)=5.
\]
故
\[
\dim U_0=5,\qquad \dim U_{\chi}=\dim U_{\bar\chi}=4.
\]
Prym 余切空间正对应于非平凡特征部分
\[
H^0(A,\Omega_A^1)\simeq U_{\chi}\oplus U_{\bar\chi},
\]
因此 \(K\)-作用在 \(H^{1,0}(A)\) 上的两条复嵌入重数均为 \(4\)，即 signature 为 \((4,4)\)。

另一方面，甲板变换群 \(C_3\) 通过保持 Prym 诱导极化的自同构作用于 \(A\)，故 Rosati 对合把 \(\sigma\) 送到 \(\sigma^{-1}\)，也即在 \(K=\QQ(\zeta_3)\) 上对应复共轭。于是 \(A\) 在该极化下确为 \(K\)-Weil 型阿贝尔八维簇。对签名 \((4,4)\) 的 Weil 型阿贝尔簇，一般理论表明
\[
W_K=\bigwedge\nolimits_K^{8}H^1(A,\QQ)
\]
是二维 \(\QQ\)-向量空间，并落在 \(H^8(A,\QQ)\cap H^{4,4}(A)\) 中。最后，同源 \(A\sim_{\QQ}J(Z)^2\) 在有理上同调上诱导 Hodge 结构同构，故 \(W_K\) 传递到 \(H^8(J(Z)^2,\QQ)\) 中相应的 \((4,4)\)-部分。证毕。
\end{proof}

\begin{conclusion}[Burnside--Kani--Rosen 同源关系诱导的局部 \texorpdfstring{$\zeta$}{zeta} 因子分解与点数恒等式]\label{con:xi-time-part9n-burnside-kani-rosen-local-zeta-pointcount}
设 \(v\) 为一处使
\[
X,\ H,\ C_6,\ Z,\ C_F,\ Y,\ X/C_2,\ X/C_3
\]
及阿贝尔簇 \(P\) 同时良约化的有限素位，记残域大小为 \(q_v\)。对任意光滑射影曲线 \(C/\FF_{q_v}\)，定义
\[
Z(C/\FF_{q_v},T)=\frac{P_C(v,T)}{(1-T)(1-q_vT)},
\]
其中 \(P_C(v,T)\) 为 Frobenius 在 \(H^1_{\text{ét}}(C_{\overline{\FF}_{q_v}},\QQ_\ell)\) 上的特征多项式；并记
\[
P_P(v,T):=
\det\!\Bigl(1-T\,\mathrm{Frob}_v\mid
H^1_{\text{ét}}(P_{\overline{\FF}_{q_v}},\QQ_\ell)\Bigr).
\]
则有
\[
P_{C_6}(v,T)=P_H(v,T)\,P_Z(v,T)^2,
\]
\[
P_{X/C_2}(v,T)=P_{C_F}(v,T)^2\,P_Y(v,T),
\]
\[
P_{C_6}(v,T)\,P_{X/C_3}(v,T)^2
=
P_H(v,T)^3\,P_{C_F}(v,T)^2\,P_Y(v,T)^2,
\]
以及
\[
P_X(v,T)=P_H(v,T)\,P_Z(v,T)^2\,P_{C_F}(v,T)^3\,P_P(v,T)^3.
\]

等价地，对任意 \(n\ge 1\)，若记
\[
N_C(v,n):=\#C(\FF_{q_v^n}),
\qquad
a_C(v,n):=q_v^n+1-N_C(v,n),
\]
则有精确恒等式
\[
a_{C_6}(v,n)=a_H(v,n)+2a_Z(v,n),
\]
\[
a_{X/C_2}(v,n)=2a_{C_F}(v,n)+a_Y(v,n),
\]
\[
a_{C_6}(v,n)+2a_{X/C_3}(v,n)
=
3a_H(v,n)+2a_{C_F}(v,n)+2a_Y(v,n).
\]
\end{conclusion}
\begin{proof}
由定理 \ref{thm:killo-s4-burnside-kani-rosen-prym-square}，
\[
J(C_6)\sim_{\QQ}J(H)\times J(Z)^2,
\]
\[
J(X/C_2)\sim_{\QQ}J(C_F)^2\times J(Y),
\]
\[
J(C_6)\times J(X/C_3)^2
\sim_{\QQ}
J(H)^3\times J(C_F)^2\times J(Y)^2.
\]
再由定理 \ref{thm:killo-s4-genus49-jacobian-computable-isogeny}，
\[
J(X)\sim_{\QQ}J(H)\times J(Z)^2\times J(C_F)^3\times P^3.
\]
在共同良约化处，同源阿贝尔簇的 \(\ell\)-进 \(H^1\) 具有相同的 Frobenius 特征多项式，而直积对应特征多项式相乘，故立即得到四条 \(P\)-因子恒等式。

对任意曲线 \(C/\FF_{q_v}\)，标准展开给出
\[
\log P_C(v,T)=-\sum_{n\ge 1}\frac{a_C(v,n)}{n}T^n.
\]
将前三条 \(P\)-因子恒等式取对数并比较 \(T^n\) 系数，即得三条点数恒等式。证毕。
\end{proof}

\begin{conclusion}[九维 Prym 因子的 Frobenius 迹整除律]\label{con:xi-time-part9n-prym9-frobenius-trace-divisibility}
对任意共同良约化素位 \(v\) 及任意 \(n\ge 1\)，有严格公式
\[
\operatorname{Tr}\!\Bigl(
\mathrm{Frob}_v^n\mid
H^1_{\text{ét}}(P_{\overline{\FF}_{q_v}},\QQ_\ell)
\Bigr)
=
\frac{
a_X(v,n)-a_H(v,n)-2a_Z(v,n)-3a_{C_F}(v,n)
}{3}.
\]
特别地，
\[
a_X(v,n)-a_H(v,n)-2a_Z(v,n)-3a_{C_F}(v,n)\equiv 0\pmod 3.
\]
\end{conclusion}
\begin{proof}
由定理 \ref{thm:killo-s4-genus49-jacobian-computable-isogeny}，
\[
J(X)\sim_{\QQ}J(H)\times J(Z)^2\times J(C_F)^3\times P^3.
\]
因此在共同良约化处，
\[
\operatorname{Tr}\!\Bigl(
\mathrm{Frob}_v^n\mid H^1_{\text{ét}}(J(X)_{\overline{\FF}_{q_v}},\QQ_\ell)
\Bigr)
=
a_H(v,n)+2a_Z(v,n)+3a_{C_F}(v,n)
\]
\[
\qquad\qquad
+3\,\operatorname{Tr}\!\Bigl(
\mathrm{Frob}_v^n\mid H^1_{\text{ét}}(P_{\overline{\FF}_{q_v}},\QQ_\ell)
\Bigr).
\]
又对曲线 Jacobian 有
\[
H^1_{\text{ét}}(J(X)_{\overline{\FF}_{q_v}},\QQ_\ell)
\cong
H^1_{\text{ét}}(X_{\overline{\FF}_{q_v}},\QQ_\ell),
\]
故左端正是 \(a_X(v,n)\)。移项并除以 \(3\) 即得所示公式。右端既等于实际的 Frobenius 迹，因而必为整数，从而得到模 \(3\) 整除律。证毕。
\end{proof}

\begin{conclusion}[属 \(49\) Jacobian 上由 \texorpdfstring{$\QQ(\zeta_3)$}{Q(zeta3)}-Weil 型强制的余维 \texorpdfstring{$4$}{4} Hodge 类]\label{con:xi-time-part9n-jx-codim4-weil-hodge-classes}
在 \(H^8(J(X),\QQ)\) 中存在自然二维 \(\QQ\)-子空间
\[
W_{K,X}\subset H^8(J(X),\QQ)\cap H^{4,4}(J(X)),
\qquad K=\QQ(\zeta_3),
\]
它由结论 \ref{con:xi-time-part9n-prym-c6h-weil-type-signature} 中 \(A=\operatorname{Prym}(C_6/H)\) 的 Weil classes 经同源分解
\[
J(X)\sim_{\QQ}J(H)\times J(Z)^2\times J(C_F)^3\times P^3
\]
传递而来。换言之，属 \(49\) 的 Jacobian \(J(X)\) 含有一组由无分歧循环三重覆盖 \(C_6\to H\) 强制产生的余维 \(4\) Hodge 类。
\end{conclusion}
\begin{proof}
由结论 \ref{con:xi-time-part9n-prym-c6h-weil-type-signature}，存在二维 \(\QQ\)-向量空间
\[
W_K\subset H^8(A,\QQ)\cap H^{4,4}(A),
\qquad A=\operatorname{Prym}(C_6/H).
\]
同一结论与定理 \ref{thm:killo-s4-burnside-kani-rosen-prym-square} 给出
\[
A\sim_{\QQ}J(Z)^2.
\]
故 \(W_K\) 可视为
\[
H^8\!\bigl(J(Z)^2,\QQ\bigr)\cap H^{4,4}\!\bigl(J(Z)^2\bigr)
\]
中的二维 Hodge 子结构。

另一方面，定理 \ref{thm:killo-s4-genus49-jacobian-computable-isogeny} 给出
\[
J(X)\sim_{\QQ}J(H)\times J(Z)^2\times J(C_F)^3\times P^3.
\]
该同源分解在有理上同调上诱导 Hodge 结构同构，因此
\[
H^8\!\bigl(J(Z)^2,\QQ\bigr)
\]
自然识别为 \(H^8(J(X),\QQ)\) 的一个有理 Hodge 直和因子。将上述 \(W_K\) 沿此同构拉回，即得
\[
W_{K,X}\subset H^8(J(X),\QQ)\cap H^{4,4}(J(X)),
\]
且 \(\dim_{\QQ}W_{K,X}=2\)。证毕。
\end{proof}
